% =========================================================================
% p1_sec08_conclusion.tex — Part I conclusion, structural nomenclature.
% =========================================================================
\section{Conclusion and Outlook of Part~I}
\label{sec:conclusion}

This Part~I of the SENI-142 study contributed the open structural layer of a measured national transmission benchmark: an open-access, georeferenced model of the Dominican bulk grid --- 142 nodal buses (115 bulk buses and 27 retained 69~kV boundary ports derived from 112 substation estates), 135 lines, 58 transformers, and ten territorial microgrids under Law 368-22 --- released with complete parameters, provenance, and validation harness via GitHub. The HBCF diakoptic partition into four macro-subdomains admits an exact Schur condensation with complexity $O(\sum_k N_k^3 + N_B^3)$ and a measured algebraic residual $\delta = 2.788 \times 10^{-14}\ \mathrm{pu}$. A four-phase fail-closed digital-twin lifecycle with detached SHA-512 custody guarantees determinism: the data verifier passes 41/41 checks, and the published mutation suite is detected 12/12 under canonical corruptions with a passing control. The resilience layer formalizes the operator's 215-contingency catalog under an explicit two-layer reserve doctrine and a four-clause static PASS battery.

Two research routes follow at the structural level: (i) \textbf{physics-informed neural surrogates} trained on the released hourly horizon with exact AC solvers as falsification oracles \cite{misyris2020}; and (ii) \textbf{grid-forming EMT co-simulation} embedding converter dynamics \cite{rosso2021}, \cite{ieee2800} to investigate frequency formation. The companion Part~II~\cite{seni142part2} completes the study with the multi-vector flexibility layer, the 87,360-state campaign, severe common-mode contingencies, and insular policy implications.
